% !TeX root = ../教材题库_全部条目_转写工作稿.tex
% Source: XJ-B1-C06-CA-P224-Y0720
% 第6章 / 6.3 统计图表 / 教材 p.224、225、226 / PDF 第231、232、233页
% OCR未校对：公式与文字均可能有误；未完成数学TeX转写；禁止直接发给学生
% 原题及相邻原解截图（核对用）：../题目原图/XJ-B1-C06-CA-P224-Y0720-01.webp, ../题目原图/XJ-B1-C06-CA-P224-Y0720-02.webp, ../题目原图/XJ-B1-C06-CA-P224-Y0720-03.webp
% 本文件只在转写完成后作为题干使用；OCR 条目尚不是可靠题干。
\begin{bankitem}{6.3 案例 案例1}{第6章 / 6.3 统计图表 / 教材 p.224、225、226 / PDF 第231、232、233页}{XJ-B1-C06-CA-P224-Y0720}{OCR未校对：公式与文字均可能有误；未完成数学TeX转写；禁止直接发给学生}
\begin{ocrdraft}
[教材p.224 / PDF第231页]
\par
案例1为研究不同类型饮料的市场销售情况，一家市场调查公司对随机抽取
\par
的一家超市进行调查．
\par
表6\_1是调查员随机观察50名顾客购买饮料类型的记录：
\par
表6一1
\par
顾客性别饮料类型
\par
顾客性别
\par
饮料类型
\par
顾客性别
\par
饮料类型
\par
顾客性别
\par
饮料类型
\par
男
\par
碳酸饮料
\par
女
\par
矿泉水
\par
女
\par
碳酸饮料
\par
女
\par
其他
\par
女
\par
茶饮料
\par
男
\par
其他
\par
女
\par
茶饮料
\par
男
\par
碳酸饮料
\par
男
\par
矿泉水
\par
男
\par
碳酸饮料
\par
男
\par
茶饮料
\par
女
\par
果汁
\par
女
\par
矿泉水
\par
女
\par
茶饮料
\par
男
\par
碳酸饮料
\par
男
\par
矿泉水
\par
女
\par
碳酸饮料
\par
女
\par
碳酸饮料
\par
女
\par
碳酸饮料
\par
男
\par
其他
\par
男
\par
矿泉水
\par
女
\par
其他
\par
女
\par
茶饮料
\par
女
\par
碳酸饮料
\par
男
\par
碳酸饮料
\par
男
\par
矿泉水
\par
男
\par
矿泉水
\par
女
\par
其他
\par
女
\par
茶饮料
\par
女
\par
碳酸饮料
\par
女
\par
茶饮料
\par
男
\par
果汁
\par
女
\par
男
\par
茶饮料
\par
男
\par
碳酸饮料
\par
女
\par
茶饮料
\par
男
\par
碳酸饮料
\par
男
\par
其他
\par
女
\par
矿泉水
\par
女
\par
果汁
\par
男
\par
碳酸饮料
\par
男
\par
碳酸饮料
\par
男
\par
果汁
\par
女
\par
其他
\par
女
\par
男
\par
茶饮料
\par
男
\par
茶饮料
\par
男
\par
矿泉水
\par
女
\par
矿泉水
\par
女
\par
其他
\par
（1)试根据上述抽样信息，绘制频数分布表．
\par
（2）试用扇形统计图、条形统计图来表示顾客购买不同类型饮料的情况．
\par
[教材p.225 / PDF第232页]
\par
解
\par
（D将样本数据按类型分类整理，制成如下频数分布表：
\par
购买数量（频数）
\par
饮料类型
\par
合计
\par
男
\par
女
\par
2
\par
4
\par
6
\par
矿泉水
\par
6
\par
4
\par
10
\par
茶饮料
\par
4
\par
7
\par
碳酸饮料
\par
9
\par
6
\par
3
\par
5
\par
8
\par
24
\par
26
\par
50
\par
从频数分布表可以看出，样本中购买碳酸饮料的顾客最多，购买矿泉水和茶饮
\par
料的顾客较多，而购买果汁的顾客最少．
\par
（2）为了直观地看出顾客购买饮料类型的情况，可以借助扇形统计图来予以呈
\par
现，如图6，3一1.
\par
在扇形统计图中，我
\par
12％
\par
们用圆形及圆内扇形的角
\par
30％
\par
0
\par
果汁
\par
20％
\par
0矿泉水
\par
度来表示一个样本（或总
\par
0茶饮料
\par
体）中各组成部分的数据
\par
16％
\par
0碳酸饮料
\par
占整体数据的比例．这对
\par
22％
\par
0
\par
其他
\par
于研究结构性问題十分
\par
有用，
\par
图6．3一1
\par
为了综合体现顾客性别对选购饮料类型的差异，可以借助复式扇形统计图以及
\par
复式条形统计图来予以呈现，如图6．3一2．
\par
10
\par
顾客性别
\par
9
\par
男
\par
8
\par
。果汁
\par
7
\par
6
\par
过矿泉水
\par
5
\par
到茶饮料
\par
4
\par
3
\par
。碳酸饮料
\par
2
\par
。其他
\par
0
\par
果汁
\par
矿泉水茶饮料碳酸饮料其他
\par
（2）
\par
图6．3一2
\par
根据图6．3一2，可以分析出该超市男、女顾客对饮料类型的喜爱程度．
\par
例如，
\par
女性顾客更多购买茶饮料，而男性顾客更多购买碳酸饮料，等等．
\par
[教材p.226 / PDF第233页]
\par
若数据是在不同时间上取得的，则可以借助折线统计图清晰地反映数据的发展
\par
变化趋势．
\end{ocrdraft}
\end{bankitem}
